\documentclass[twocolumn,a4paper]{article}
\usepackage[a4]{draft}

\usepackage{color}
% \usepackage{txfonts}
\usepackage{cite}
\usepackage{amsmath,amssymb,amsfonts}
\usepackage{algorithmic}
\usepackage{graphicx}
\usepackage{mediabb}
\usepackage{fancybox}

% 日本語設定 (LuaLaTeX)
\usepackage{luatexja}
% \usepackage[haranoaji]{luatexja-preset} % タイムアウト回避のためコメントアウト中

\begin{document}
%date not printed
\date{}

%make title
\title{\Large\bf
鳥獣害対策のための複数台のカメラを用いた低コスト通信型監視システムの開発
\\~\\
\large\bf
Development of a Low-Cost Communication-Based Monitoring System Using Multiple Cameras for Wildlife Damage Control}


%for two authors (Template style)
\author{\normalsize
 \begin{tabular}[t]{c@{\extracolsep{8em}}c}
  \large 著者名 (Name)& \large 共著者名 (Co-author) \\
  \\
   所属 (Affiliation) & 所属 (Affiliation) \\
   City, Country & City, Country \\
   email@example.com & email@example.com
\end{tabular}}

\maketitle

\thispagestyle{empty}

{\small\bf Abstract---
農作物への鳥獣被害は深刻な社会問題となっており、その対策としてIoT技術を用いた安価で効率的な監視システムが求められている。本研究では、鳥獣害対策を目的とし、複数台の安価な無線カメラ端末（ESP32）とデータを集約するエッジサーバー（Raspberry Pi）を用いた低コストな通信型監視システムを開発した。本システムは、エッジ側でのYOLOv8nを用いた一次フィルタリングと「3枚バースト撮影・多数決判定」により、PIRセンサー特有の誤検知を大幅に削減し、真に必要な検知情報のみを管理者に通知する。これにより、通信コストの抑制と、被害発生時の迅速な対応支援を実現する。
}

\section{はじめに (Introduction)}
近年、シカやイノシシなどの野生動物による農作物被害や生活圏への出没が深刻化しており、鳥獣害対策は喫緊の課題である。被害防除のためには、動物の出没状況を正確かつリアルタイムに把握することが重要である。カメラトラップはその有効な手段の一つであるが、従来のPIR（受動型赤外線）センサーを用いた安価なカメラは、風による植物の揺れなどを誤検知することが多く、大量の「空撮影」データにより確認作業が煩雑になるという問題がある\cite{beery2019iwildcam}。また、セルラー通信を用いて画像をリアルタイム送信する場合、誤検知画像の送信は通信コストの浪費につながる。

一方で、広範囲な農地や里山を監視するためには、多数のカメラを設置する必要があり、システム全体の低コスト化が不可欠である。そこで本研究では、鳥獣害対策への適用を念頭に、複数台の安価なカメラ端末とそれらを統括するエッジサーバーから成る、低コストかつ高効率な通信型監視システムを提案する。具体的には、Raspberry Piを用いたエッジコンピューティングにより、現場で画像の一次選別を行う（図\ref{fig:system_overview}参照）。

\section{関連研究 (Related Work)}
野生動物モニタリングの分野では、従来よりPIRセンサーを用いたトレイルカメラが広く用いられてきた。しかし、Beeryら\cite{beery2019iwildcam}が指摘するように、PIRセンサーは温度変化や風による揺れに対しても反応するため、撮影画像の大半（時には90\%以上）が動物を含まない「空撮影（empty images）」となることが知られている。この問題に対し、Norouzzadehらによるディープラーニングを用いた大規模画像分類の試みがあるが、これらは主に撮影後のオフライン解析に焦点を当てている。

IoT技術の発展に伴い、エッジデバイス上でリアルタイムに動物検知を行う「スマートカメラトラップ」の研究も進んでいる。例えば、Raspberry Pi等のシングルボードコンピュータにMobileNetやYOLO等の軽量モデルを搭載する手法である\cite{redmon2016yolo}。しかし、常時映像を解析する手法は消費電力が大きく、商用電源のない山間部での長期運用には適さない。

本研究の独自性は、(1) ESP32による超低消費電力な待機・撮影機能と、(2) Raspberry Piによる高精度な推論機能を物理的に分離し、Wi-Fiで連携させる階層型アーキテクチャにある。さらに、PIRの誤検知特性を考慮した「3枚バースト撮影・多数決フィルタリング」を導入することで、単一の画像解析のみに頼るシステムよりもロバストな検知を実現している点が最大の特徴である。

\section{システム設計 (System Design)}
提案システムは、以下の3つのレイヤーで構成される「階層型処理アーキテクチャ」を採用している。

\subsection{エッジデバイス (Sensing Node)}
撮影端末として、ESP32-S3マイコンを採用したカメラデバイスを使用する。
\begin{enumerate}
    \item \textbf{省電力動作}: 通常時はDeep Sleepモードで待機し、PIRセンサー検知時のみ起動する。
    \item \textbf{バースト撮影}: 誤検知のリスクを考慮し、一度の検知で3枚の画像を連続撮影する。この「3枚1サイクル」を処理単位とする。
    \item \textbf{メタデータ付与}: GPSモジュールにより取得した正確な位置情報と時刻を付与する。
\end{enumerate}

\subsection{エッジサーバー (Edge Gateway)}
複数の撮影端末からのデータを集約するゲートウェイとして、Raspberry Piを使用する。
\begin{itemize}
    \item \textbf{一次フィルタリング (YOLOv8n)}: 受信した画像に対し、軽量な物体検出モデルであるYOLOv8nを実行する。Raspberry Pi 4上での推論速度と精度のバランスを考慮し、Nanoモデルを採用した。
    \item \textbf{多数決ロジック}: 3枚1サイクルの画像のうち、2枚以上で動物クラス（bird, cat, dog等）が検出された場合のみ、そのサイクルを「有効」と判定する。単一画像のみの誤検出を排除し、信頼性を向上させる。
\end{itemize}

\subsection{外部サーバー (Analysis Server)}
最終的なデータ蓄積と解析を行うクラウドまたはオンプレミスサーバーである。
\begin{itemize}
    \item \textbf{詳細解析 (MegaDetector)}: 計算資源の豊富な環境でMegaDetector v5を使用し、動物種の特定とカウントを行う。
    \item \textbf{通知}: 解析結果に基づき、検知ボックス（Bounding Box）を描画した画像を管理者へメール通知する。
\end{itemize}

% 図のプレースホルダー
\begin{figure}[t]
  \centering
  % \includegraphics[width=8cm]{fig/system_overview.pdf}
  \fbox{\parbox{0.9\linewidth}{
      \centering
      \vspace{2cm}
      System Overview Diagram\\
      (Device -> Edge -> Cloud)
      \vspace{1cm}
  }}
  \caption{提案システムの概要。エッジサーバーでの一次フィルタリングにより、外部サーバーへの送信データ量を削減する。}
  \label{fig:system_overview}
\end{figure}

\section{実装 (Implementation)}
\subsection{ハードウェア構成}
撮影ノードにはSeeed Studio XIAO ESP32S3 Senseを使用し、外部アンテナによりWi-Fi通信距離を確保した。ゲートウェイにはRaspberry Pi 4 Model B (4GB RAM) を使用した。

\subsection{ソフトウェア実装}
エッジデバイスのファームウェアはArduino IDE (C++) で開発し、HTTP/multipart形式で画像を転送する。エッジサーバーおよび外部サーバーはPythonで記述され、Flaskを用いたREST APIサーバーとして機能する。

\section{評価 (Evaluation)}
本提案システムの有効性を検証するため、予備実験を行った。（ここに具体的な実験結果を記述）

\section{結論 (Conclusion)}
本稿では、階層的なフィルタリング機能を持つ野生動物モニタリングシステムを提案した。エッジ側での予備選別により、不要なデータ送信を抑制し、効率的なモニタリングが可能となる。

\section*{\sc Acknowledgements}
本研究の一部は、[助成金名など]の支援を受けたものである。

\footnotesize

\bibliographystyle{myjunsrt}
\bibliography{references}

\end{document}
